\section{Introduction}
\label{sec:intro}

The models used throughout this thesis describe populations small enough that chance
is part of the mechanism. A handful of bacteria deposited in a mammalian host, a few
viral particles inside a single cell: in each case the count lies far from the regime
in which deterministic rate equations are informative. The questions that matter ---
will the population establish at all, and if it survives, how large will it be? ---
are questions about fluctuations rather than about a single mean trajectory. This
chapter assembles the mathematical apparatus needed to answer them, and to prepare the
birth--death--catastrophe and multi-type models of later chapters.

Two threads organise what follows.

The first is \emph{conditioning on survival}. A subcritical branching process dies out
with probability one, so its unconditional mean decays geometrically and conveys
little about realisations that remain alive. Looking only at those realisations that
are still positive at generation $n$ yields a different picture: the conditional
population settles to a fixed distribution, the quasi-stationary (or Yaglom) limit,
whose mean is a constant. For the binary Galton--Watson process with replication
probability $p<\tfrac12$ that constant is $1/A(p)$, where
\begin{equation*}
A(p) = \lim_{n\to\infty}\frac{S_n}{(2p)^n},
\end{equation*}
and $S_n$ is the probability of survival to generation $n$. Much of the chapter is
concerned with $A(p)$: what it is, how to compute it, and why --- despite sustained
effort --- no elementary closed form has been found.

The second thread is the \emph{near-critical regime}. As $p$ approaches $\tfrac12$ from
below the process takes ever longer to die, the quasi-stationary mean $1/A(p)$
diverges, and every numerical scheme that relies on iterating to stationarity becomes
impractical. Just above $\tfrac12$ the initial count of replicative individuals weighs
heavily on the chance that the process survives at all. Behaviour on either side of
$p=\tfrac12$ is what makes small inocula interesting, and what makes them awkward to
compute with.

The chapter is ordered as follows. \Cref{sec:prelim} records the short continuous-time
toolkit used later: exponential races, continuous-time Markov chains and their
generators, probability generating functions, and first-step analysis.
\Cref{sec:gw} introduces the binary Galton--Watson process and establishes the two
facts on which the rest rests: extinction is certain at and below criticality, and at
criticality the expected lifetime is nevertheless infinite. \Cref{sec:qsd} derives the
limiting conditional mean and variance and so introduces $A(p)$. \Cref{sec:small}
treats founding cohorts of size $k$, discrete birth--death--catastrophe survival
formulae, and the continuous-time birth--death constant $A_{\mathrm{c}}(p)$, for which a
closed form is available. \Cref{sec:Ap} is devoted to $A(p)$ itself: an infinite-product
representation, an exact series with two-sided bounds, the near-critical asymptotic
$A(p)\sim 2(1-2p)$, a comparison of discrete and continuous constants, symbolic-regression
searches, and the Koenigs-linearisation argument that explains the absence of an
elementary formula. \Cref{sec:moc} develops generating-function partial differential
equations and the method of characteristics for the absorption models used later in the
thesis.

Results that are standard are presented as such. Places where the argument is heuristic,
or where numerics carry the load, are marked explicitly. Where a claim can be checked
numerically, it has been.
